\documentclass[12pt]{article}
\usepackage[utf8]{inputenc}
\usepackage{amsmath}
\usepackage{amssymb}
\usepackage{graphicx}
\usepackage{booktabs}
\usepackage{listings}
\usepackage{xcolor}
\usepackage{hyperref}
\usepackage{geometry}
\geometry{margin=1in}

\lstset{
    language=Python,
    basicstyle=\ttfamily\small,
    keywordstyle=\color{blue}\bfseries,
    commentstyle=\color{green!60!black},
    stringstyle=\color{red},
    numbers=left,
    numberstyle=\tiny\color{gray},
    stepnumber=1,
    numbersep=5pt,
    backgroundcolor=\color{gray!10},
    frame=single,
    breaklines=true,
    breakatwhitespace=false,
    tabsize=4,
    showstringspaces=false
}

\title{Environmental Mitigation Costs Calculation Documentation}
\author{}
\date{}

\begin{document}

\maketitle

\section{Introduction}

The \texttt{environmental\_mitigation.py} script calculates environmental mitigation costs for a transmission line project. It computes base construction/restoration costs and wetland/habitat credit purchases across different construction types and terrain types, applying uplift factors for effective acreage.

\section{Method Overview}

\subsection{Purpose}

This script calculates environmental mitigation costs by:
\begin{enumerate}
    \item Calculating ROW acres by terrain type
    \item Applying uplift factor to determine effective acres (TCE - Total Cost Equivalent)
    \item Computing base mitigation/restoration costs per effective acre
    \item Calculating wetland credit purchases for impacted wetland acres
    \item Calculating habitat credit purchases for impacted natural habitat acres
    \item Applying AFUDC capitalization for regulatory perspective
    \item Computing present values for societal perspective
\end{enumerate}

\subsection{Calculation Approach}

The environmental mitigation cost calculation follows this structure:

\begin{align}
\text{Terrain Acres} &= \frac{\text{Terrain Miles} \times 5280 \times \text{ROW Width (feet)}}{43560} \\
\text{Effective Acres} &= \text{Base Acres} \times \text{Uplift Factor} \\
\text{Base Cost} &= \sum_{\text{terrains}} \text{Cost/Acre} \times \text{Effective Acres} \\
\text{Wetland Credits} &= \text{Wetland Cost/Acre} \times \text{Ratio} \times \text{Wetland Impacted Acres} \\
\text{Habitat Credits} &= \text{Habitat Cost/Acre} \times \text{Ratio} \times \text{Habitat Impacted Acres} \\
\text{Total} &= \text{Base Cost} + \text{Wetland Credits} + \text{Habitat Credits}
\end{align}

The uplift factor accounts for additional mitigation requirements beyond the physical ROW footprint (e.g., buffer zones, restoration areas).

\subsection{Inputs and Outputs}

\textbf{Inputs:}
\begin{itemize}
    \item Project technical details (construction type, AC/DC, capacity, etc.)
    \item Physical details (terrain distribution)
    \item ROW width for project category
    \item Environmental mitigation parameters (base costs, credit costs, credit ratios, uplift factor)
    \item Financing parameters (WACC, inflation, AFUDC rate)
    \item Timeline (delay years, construction years)
\end{itemize}

\textbf{Outputs:}
\begin{itemize}
    \item Nominal costs (base mitigation, wetland credits, habitat credits)
    \item AFUDC capitalized costs (regulatory perspective)
    \item Present value costs (societal perspective)
    \item Acreage metrics (base acres, effective acres, impacted acres)
\end{itemize}

\subsection{Integration with Other Scripts}

This script integrates with:
\begin{itemize}
    \item \texttt{yaml\_loaders.py} - Loads project configuration and environmental mitigation parameters (see \texttt{utilities\_documentation.tex})
    \item \texttt{financial\_utils.py} - Financial calculations (see \texttt{utilities\_documentation.tex})
    \item \texttt{csv\_output\_manager.py} - Writes results to CSV (see \texttt{csv\_output\_manager\_documentation.tex})
\end{itemize}

\section{Key Functions}

\subsection{\texttt{calculate\_environmental\_mitigation\_costs(em\_yaml, category, terrain\_miles\_dict, row\_width\_feet)}}

Calculates environmental mitigation costs including base mitigation and credit purchases.

\textbf{Parameters:}
\begin{itemize}
    \item \texttt{em\_yaml} (dict): Loaded environmental mitigation YAML data
    \item \texttt{category} (str): Project category identifier
    \item \texttt{terrain\_miles\_dict} (dict): Dictionary mapping terrain type to miles
    \item \texttt{row\_width\_feet} (float): Right-of-way width in feet
\end{itemize}

\textbf{Returns:}
\begin{itemize}
    \item Dictionary containing:
    \begin{itemize}
        \item \texttt{base\_cost} (float): Base mitigation/restoration cost
        \item \texttt{wetlands\_credits} (float): Wetland credit purchase cost
        \item \texttt{habitat\_credits} (float): Habitat credit purchase cost
        \item \texttt{total\_credits} (float): Sum of wetland and habitat credits
        \item \texttt{total} (float): Total mitigation cost
        \item \texttt{uplift\_factor\_applied} (float): Uplift factor used
        \item \texttt{total\_base\_acres} (float): Base ROW acres
        \item \texttt{total\_effective\_acres} (float): Effective acres with uplift
        \item \texttt{wetland\_impacted\_acres} (float): Wetland acres impacted
        \item \texttt{habitat\_impacted\_acres} (float): Habitat acres impacted
    \end{itemize}
\end{itemize}

\textbf{Key Logic:}
\begin{lstlisting}
# Calculate acres by terrain
for terrain, miles in terrain_miles_dict.items():
    if miles > 0:
        terrain_acres = (miles * 5280 * row_width_feet) / 43560
        effective_acres = terrain_acres * uplift_factor
        cost_per_acre = base_costs[construction_type][terrain]
        base_cost += cost_per_acre * effective_acres

# Calculate credit purchases
wetland_impacted_acres = acres_by_terrain["wetland"] * uplift_factor
wetlands_credits = (wetland_cost_per_acre * wetland_ratio 
                     * wetland_impacted_acres)

habitat_impacted_acres = sum(acres_by_terrain[t] for t in habitat_terrains) 
                         * uplift_factor
habitat_credits = (habitat_cost_per_acre * habitat_ratio 
                   * habitat_impacted_acres)
\end{lstlisting}

\subsection{\texttt{main()}}

Main execution function that orchestrates the environmental mitigation cost calculation and output.

\textbf{Workflow:}
\begin{enumerate}
    \item Loads project technical details and constructs category identifier
    \item Loads ROW width and terrain distribution
    \item Loads environmental mitigation parameters
    \item Calls \texttt{calculate\_environmental\_mitigation\_costs()} to compute costs
    \item Calculates AFUDC capitalized costs (base and credits separately)
    \item Calculates present values (base spread over construction, credits upfront)
    \item Displays results and writes to CSV
\end{enumerate}

\section{Example}

The following example demonstrates a typical usage of the environmental mitigation cost calculation:

\begin{lstlisting}
# Example: 500 MW Overhead AC line
# Category: "Overhead/AC/500MW/Advanced Aluminum/NA"
# 
# Inputs (from YAML):
#   - ROW width: 200 feet
#   - Terrain distribution:
#     * Forested: 40 miles
#     * Wetland: 10 miles
#     * Farmland: 50 miles
#   - Uplift factor: 1.5x
#   - Base cost: $2000/acre (overhead, forested), $3000/acre (overhead, wetland)
#   - Wetland credit: $5000/acre, ratio 2.0:1
#   - Habitat credit: $3000/acre, ratio 1.5:1
#   - Construction: 2 years
#
# Calculation:
#   Forested acres = (40 * 5280 * 200) / 43560 = 969.7 acres
#   Wetland acres = (10 * 5280 * 200) / 43560 = 242.4 acres
#   Farmland acres = (50 * 5280 * 200) / 43560 = 1212.1 acres
#
#   Effective acres (with 1.5x uplift):
#   * Forested: 969.7 * 1.5 = 1454.6 acres
#   * Wetland: 242.4 * 1.5 = 363.6 acres
#
#   Base cost = (1454.6 * $2000) + (363.6 * $3000) = $4.4M
#   Wetland credits = $5000 * 2.0 * 363.6 = $3.6M
#   Habitat credits = $3000 * 1.5 * 1454.6 = $6.5M
#   Total = $14.5M
#
# Output:
#   - Nominal cost: $14.5M
#   - AFUDC capitalized: $15.2M (if AFUDC enabled)
#   - Present value: $13.8M (base spread over 2 years, credits upfront)
\end{lstlisting}

\textbf{Integration Example:}

\begin{lstlisting}
# When called from ctcc.py:
csv_manager = CTCCOutputManager()
csv_results = {
    "total_nominal": results["total"],
    "total_afudc": total_capitalized if apply_afudc else 0,
    "total_pv": total_pv,
    "base_cost_nominal": results["base_cost"],
    "credits_nominal": results["total_credits"],
}
csv_manager.add_environmental_mitigation(csv_results)
csv_manager.write_batch_summary()
\end{lstlisting}

\section{Key Implementation Details}

\subsection{Uplift Factor}

The uplift factor (TCE - Total Cost Equivalent) accounts for mitigation requirements beyond the physical ROW footprint. For example, a 1.5x uplift means that for every acre of ROW, 1.5 acres of mitigation is required (including buffer zones, restoration areas, etc.).

\subsection{Credit Ratios}

Credit ratios represent the mitigation requirement. For example, a 2.0:1 ratio for wetlands means that for every acre of wetland impacted, 2 acres of wetland credits must be purchased.

\subsection{Habitat Terrain Classification}

The script identifies habitat-impacted acres by summing natural terrain types: forested, scrubbed\_flat, desert\_barren, and rolling\_hills. This is a conservative approach that captures all potential habitat impacts.

\subsection{Present Value Timing}

\begin{itemize}
    \item \textbf{Base mitigation costs}: Spread evenly over construction period (annual payments)
    \item \textbf{Credit costs}: One-time payment at start of construction (end of delay period)
\end{itemize}

\subsection{AFUDC Treatment}

Base mitigation and credit costs are treated separately for AFUDC calculation, as they may have different timing patterns during construction.

\end{document}

